\bookchapter{第四十七章：クイックリファレンスとリソース}

\chapterintro{この章は、これ1つで使えるクイックリファレンスです。原則をうろ覚えのときや、すばやく判断したいときは、直接ここを開いてください。「ポケットカード」として手元に置くのがおすすめです。5大原則の早見表、ツール選択の決定木、よくある質問、学習リソース——どれも飛ばさずどうぞ。}

\booksection{思考モデルの早見カード}

\begin{center}
\begin{tabularx}{\linewidth}{|X|X|X|}
\hline
\textbf{原則} & \textbf{ひとことで} & \textbf{操作のポイント} \tabularnewline
\hline
Think Before Coding & まず考えてから頼む & 3つの質問に答える：何をする？結果はどんな形？制約は何？ \tabularnewline
Simplicity First & 最小限のコードで問題を解く & 1回に1機能だけ。不要な複雑さは切り捨てる \tabularnewline
Surgical Changes & ピンポイントで直し、境界を越えない & 修正範囲を明確に指定。「ついでに直しておきました」は断る \tabularnewline
Goal-Driven Execution & 「何をするか」を言い、「どうやるか」は言わない & 手順ではなく目標を渡し、道筋はAIに任せる \tabularnewline
Iterate \& Verify & 生成検証修正のループ & 3層検証：機能、境界、アーキテクチャ \tabularnewline
\hline
\end{tabularx}
\end{center}

\begin{center}\small\color{gray}5大思考原則の早見カード\end{center}

\booksection{ツール選択の決定木}

AIプログラミングツールを選ぶときは、次の判断経路に沿って選びましょう。注意：「どれがいちばん優れているか」のランキングではなく、「どれが自分にいちばん合うか」の判断フレームワークです。

何をしますか？

ゼロベース、まず体験したい クラウドBuilder (v0/Bolt/Replit)を選ぶ

日常の開発、コード補完が必要 IDE内蔵AI (Cursor/Trae)を選ぶ

大型プロジェクト、一括操作が必要 ターミナルAgent (Claude Code)を選ぶ

プロトタイプで素早くアイデアを検証 クラウドBuilderかTrae SOLOを選ぶ

\begin{center}\small\color{gray}図6：AIプログラミングツール選択の決定木\end{center}

\booksection{よくある質問（FAQ）}

Q1：プログラミングがまったくできませんが、AIプログラミングは使えますか？

使えます。2026年のAIプログラミングツールは十分に強力で、自然言語で要件を伝えるだけで、AIがコードを生成してくれます。ただし2つアドバイスを。1つ目は、クラウドのBuilderツール（v0やBoltなど）から始めること。開発環境の設定がいりません。2つ目は、使いながら基礎概念を学ぶこと。文法を完璧に知る必要はありませんが、AIが書いたコードの大意を「読める」ことが、きちんと検収するうえで大切です。

Q2：どのプログラミング言語を学ぶべきですか？

何をしたいかによります。WebサイトのフロントエンドならJavaScript/TypeScript。データ処理やAIならPython。企業向けのバックエンドならJavaかGoです。でも、もっと大事なことがあります。AIプログラミングの時代には、「言語」の重要度は下がりつつあるということです。文法の細かい部分はAIが引き受けてくれます。あなたに必要なのは「プログラムの構造」の理解、つまり関数とは何か、ループとは何か、データ構造とは何かです。これらの概念は言語をまたいで共通なので、一度身につければ、どの言語に乗り換えても早いのです。

Q3：AIが書いたコードに著作権の問題はありますか？

2026年8月時点では、いまだに意見が分かれる法律上の問題です。一般的には、AIが生成したコードには明確な著作権者がいるとは見なされない司法管轄区が大多数です。とはいえ商用プロジェクトで使うなら、(1)使うAIツールの利用規約の知的財産権まわりの条項を確認する、(2)重要なコードは人間が見直して手を加え、自分の創造的な貢献を含ませる、(3)著作権で保護されている可能性のあるオープンソースコードをAIにそのままコピーさせない、ことをおすすめします。

Q4：AIはプログラマーに取って代わりますか？

当面、その心配はいりません。AIが置き換えるのは「プログラマー」という職業ではなく、「1行ずつコードを書く」という作業です。電卓が数学者を消したのではなく、面倒な手計算から解放したのと同じことです。AIプログラミングツールによって、プログラマーは「コードを書く」から「問題を考える」へと移ります。これは実際には、プログラマーの価値の上限を引き上げています。これからの競争力は「何行書けるか」ではなく、「どれだけ複雑な問題を解けるか」にあるのです。

Q5：AIの出力が正しいかどうか、どう見極めればよいですか？

3層検証法を使いましょう。第1層は機能検証——動くか？入力と出力は合っているか？これが基本中の基本です。第2層は境界検証——空の入力、極端に長い入力、異常な入力でどうなるか？ここで隠れた問題が見つかります。第3層はアーキテクチャ検証——コードの構造は妥当か？ハードコードはないか？拡張しやすいか？ここにはある程度の技術的判断力がいります。第3層まで見極められなくても、最初の2層はやりましょう。動いて、境界でも問題がなければ、初学者の大半より先を行けています。

\booksection{おすすめ学習リソース}

2026年8月時点でのおすすめ学習リソースです。インターネットのリンクは変わることがあります。リンク切れのときは、タイトルで検索すればたいてい最新版が見つかります。

\begin{center}
\begin{tabularx}{\linewidth}{|X|X|X|}
\hline
\textbf{種類} & \textbf{リソース} & \textbf{おすすめの理由} \tabularnewline
\hline
一次資料 & Andrej Karpathy氏のCLAUDE.md & 5大原則の源流。70行ながら、どの1行も珠玉 \tabularnewline
一次資料 & Anthropic公式ドキュメント（Claude Code） & ターミナルエージェントのベストプラクティスガイド \tabularnewline
チュートリアル & Vibe CodingコミュニティがまとめたSkillsコレクション & コミュニティの実践者たちがまとめた実戦テクニック集 \tabularnewline
チュートリアル & Cursor公式ドキュメントとYouTubeチャンネル & IDE内蔵AIの権威ある使い方ガイド \tabularnewline
コミュニティ & GitHubのawesome-claude-code & Claude Codeのベストプラクティスとプラグイン集 \tabularnewline
コミュニティ & Redditのr/cursorコミュニティ & 世界中のCursorユーザーによる経験の共有と議論 \tabularnewline
基礎 & MDN Web Docs (developer.mozilla.org) & フロントエンド開発で最も権威あるリファレンス \tabularnewline
基礎 & Python公式チュートリアル（docs.python.org） & Python学習の公式のスタート地点 \tabularnewline
\hline
\end{tabularx}
\end{center}

\begin{center}\small\color{gray}おすすめ学習リソース一覧\end{center}

\booksection{おわりに：道具は手に、方法論は心に}

本書の核心にある信念に立ち返りましょう。道具は3か月ごとに更新され、方法論は3年かけてようやく進化します。2026年8月の今日、AIプログラミングツールの力は驚くほど強くなりました。ひと言でWebサイトを作り、説明文だけでフルスタックアプリを生み出し、指示1つで複雑なリファクタリングを仕上げられるのです。しかし、道具が強くなるほど、方法論の重みはむしろ増します。道具の強さが、あなたの一つひとつの判断を拡大増幅するからです。

要件が明確なら、AIは効率よく実現してくれます。要件があいまいなら、AIは同じ効率で「まちがったもの」を作ります。シンプル第一を貫けば、AIはアイデアの素早い検証を助けてくれます。欲張って全部入れようとすれば、AIは技術的負債を同じ速さで積み上げます。ピンポイントで指示すれば、AIは外科手術のようにコードを最適化します。指示があいまいなら、AIは大鉈を振るうようにバグを持ち込みます。

だからこそ本書は、道具の操作より方法論の解説に多くのページを割いたのです。道具が大切でないからではありません。方法論は道具の「掛ける数」だからです。方法論が良ければ道具の価値は大きくなり、方法論がゼロなら、どんなに強力な道具も「ゼロ×いくら」のままです。

もう一人でプログラミングするのではありません。あなたは指揮官、AIはあなたの軍団です。思考の質が軍団の戦闘力を決めます。今日から方法論で自分の考えを武装し、道具で自分の能力を拡大しよう——それがAIプログラミングのすべての秘密です。

\booksection{参考文献}

\begin{enumerate}
  \item Andrej Karpathy. (2025). Vibe Coding. Social media post introducing the concept of vibe-driven programming with AI.
  \item Andrej Karpathy. (2025). CLAUDE.md - Project guidelines for AI-assisted development. GitHub repository, 70+K stars.
  \item Anthropic. (2026). Claude Code: Terminal-native AI programming. Official documentation and best practices.
  \item Cursor Team. (2026). Cursor 3: The next generation of AI-first IDE. Release notes and documentation.
  \item OpenAI. (2026). Codex CLI: Bringing AI to the terminal. Technical documentation.
  \item Model Context Protocol (MCP) Specification. (2025-2026). Open protocol for AI-tool integration.
  \item Vibe Coding Community. (2026). Community-curated skills and best practices for AI-assisted programming.
  \item Spec-Driven Development (SDD) methodology. (2026). Emerging practices in specification-first AI programming.
\end{enumerate}
